\documentclass[11pt,letterpaper]{article}
\usepackage{ad42010-suite}
\begin{document}
\DocFront{06}{Practical Architecture Modeling}{UML, C4, PlantUML, security views and automated checks}{An implementation-focused laboratory showing how contemporary diagrams and engineering evidence can populate 2022 architecture description views without treating a notation as a substitute for the required metadata.}
\InterpretationKey
\section{Model-first does not mean diagram-only}
\begin{Requirement}\FlagReq{} Architecture views must follow their governing viewpoints and contain view components; model-based components identify governing model kinds, whereas non-model components require legends. The standard prescribes neither a particular notation nor authoring tool. \Source{1, 6.6--6.8}\end{Requirement}
\begin{Guidance}Think of C4, UML and PlantUML as possible \emph{representation resources}; define their roles locally as model kinds, modelling-language specifications or tools. Map an actual diagram to a named architecture view, specify the viewpoint governing that view and record the concerns addressed. A generic PNG file is insufficient provenance by itself.\end{Guidance}
\section{Notation selection and architecture purpose}
\begin{longtable}{p{0.22\linewidth}p{0.24\linewidth}p{0.38\linewidth}}\toprule\textbf{Resource}&\textbf{Useful EoI aspect}&\textbf{Appropriate application / caveat}\\\midrule\endhead
C4 system context&Structural / environment&Shows the EoI and external actors/systems; adopt a specific context viewpoint and legend.\\
C4 container&Structural / functional&Shows deployable/application-level divisions; a ``container'' is not a 42010 view component.\\
C4 component&Structural / functional&Clarifies responsibilities inside one container; not the same as AD view components.\\
UML component diagram&Structural / interfaces&Explicit contracts and dependencies; specify stereotype and port conventions.\\
UML sequence diagram&Behavioural&Message order, alternatives and timeout behaviour for selected scenarios.\\
UML activity diagram&Functional / behavioural&Paths, conditions, concurrent work and swimlane responsibility.\\
UML state machine&Behavioural&States, guards, transitions and exceptional outcomes.\\
BPMN collaboration&Business process / organizational&Process and cross-party interactions; map tasks to technical design only where justified.\\
Data-flow / threat model&Informational / trust&Sources, sinks, boundaries and mitigations; record threat vocabulary and analytical assumptions.\\
Deployment model&Physical / structural&Allocation of runtime units to nodes and environments; map conceptual to actual configuration.\\\bottomrule\end{longtable}
\begin{Guidance}A tool may render one or several model kinds. The modelling conventions and analytical procedures belong in controlled specifications; PlantUML is not in itself an architecture viewpoint. Any implementation-specific diagram dialect needs explicit syntax/semantics assumptions.\end{Guidance}
\section{Worked example: a service architecture baseline}
\begin{Worked}\textbf{AD} \texttt{AD-PAYMENTS/1.0}. \textbf{EoI} \texttt{PAY-SVC}. \textbf{Purpose} review integration resilience and sensitive-data handling.\par
\textbf{Concerns} \texttt{CN-01}: which external parties can reach the service? \texttt{CN-02}: which calls cross a security boundary? \texttt{CN-03}: how are failures completed or retried? \texttt{CN-04}: who recovers a stalled transaction?\par
\textbf{Views} \texttt{VW-CXT-01/1} governed by \texttt{VP-CXT/1}; \texttt{VW-SEC-01/1} governed by \texttt{VP-SEC/1}; \texttt{VW-FUN-01/1} governed by \texttt{VP-FUN/1}; \texttt{VW-OPS-01/1} governed by \texttt{VP-OPS/1}.\par
\textbf{Components} \texttt{VC-CXT-01} graph, \texttt{VC-FUN-01} interaction sequence, \texttt{VC-SEC-01} trust/asset table, \texttt{VC-OPS-01} narrative runbook. The narrative receives a legend; the diagram model kinds receive their own controlled specifications.\end{Worked}
\section{Laboratory 1: EoI context with PlantUML}
\textbf{Objective}: produce an original context view component with a known EoI boundary, a typed interface set and identified external stakeholders.\par
\begin{lstlisting}[caption={Example only: examples/01-context.puml}]
@startuml
skinparam componentStyle rectangle
actor "Customer" as Customer
rectangle "Payment Service [EoI]" as PaySvc
rectangle "Identity Provider" as IdP
rectangle "Payment Provider" as PSP
Customer --> PaySvc : Submit / query payment
PaySvc --> IdP : Validate identity
PaySvc --> PSP : Authorize / capture
note bottom of PaySvc
  EoI: PAY-SVC
  AD element: VC-CXT-01
  Governing model kind: MK-CONTEXT/1
end note
@enduml
\end{lstlisting}
\textbf{Representation contract (local)}: rectangles designate named logical entities; arrows represent high-level integration relationships, not necessarily packet directions; labels are interaction purposes. \textbf{Viewpoint}: \texttt{VP-CXT/1}. \textbf{Concern coverage}: \texttt{CN-01} and part of \texttt{CN-02}.\par
\textbf{Validation exercise}: compare the external nodes with the API/service inventory. Report missing dependencies as \texttt{INC-*} and unknown interfaces as issues, not as inferred fact.\par
\section{Laboratory 2: UML-style interaction model}
\begin{lstlisting}[caption={Example only: examples/02-payment-sequence.puml}]
@startuml
actor Client
participant "Payment API" as API
queue "Durable Queue" as Queue
participant "Payment Worker" as Worker
participant "Provider" as PSP
Client -> API : POST /payments (idempotency key)
API -> Queue : publish accepted command
API --> Client : 202 Accepted + transaction ID
Queue -> Worker : deliver command
Worker -> PSP : authorize(transaction)
alt provider responds
    PSP --> Worker : success / decline
    Worker -> Queue : acknowledge + publish status
else timeout / unavailable
    Worker -> Queue : retry policy / dead-letter
end
@enduml
\end{lstlisting}
\textbf{Governance}: \texttt{VC-FUN-01/1}, \texttt{MK-SEQUENCE/1}, \texttt{VP-FUN/1}. \textbf{Framed and addressed concerns}: completion/retry \texttt{CN-03}; external provider behaviour in \texttt{CN-02}. \textbf{Known issues}: threshold and recovery semantics must be confirmed with approved requirements; this drawing does not settle them.\par
\textbf{View method (original)}: start from a named user scenario; add negative branches; trace each message to an interface or event-contract ID; review timeouts, duplicates, side effects, actor ownership and compensating actions.\par
\section{Laboratory 3: trust-boundary component}
\begin{longtable}{p{0.21\linewidth}p{0.32\linewidth}p{0.33\linewidth}}\toprule\textbf{Element}&\textbf{Trust/control interpretation}&\textbf{Evidence still needed}\\\midrule\endhead
Client to API&Untrusted-to-service boundary; authentication and authorization concerns.&Identity model; enforcement ownership; negative tests.\\
API to queue&Durable message handoff; sensitive metadata and replay concern.&Encryption/key handling; broker access policy; idempotence evidence.\\
Worker to provider&External provider dependency; integrity and exposure concerns.&Credential management; outbound control; retry safety.\\
Status update&Potential leakage and confused-deputy concern.&Object-level authorization rules; status event schema.\\\bottomrule\end{longtable}
\begin{Guidance}Treat this table as \texttt{VC-SEC-01}, a non-model narrative/table component with legend \texttt{LG-SEC-TABLE/1}. In this legend ``concern'' denotes a proposed issue to assess, not an observed vulnerability; ``evidence'' means a record to obtain. Threat classifications such as STRIDE, CWE and OWASP are optional specialist extensions external to ISO/IEC/IEEE 42010.\end{Guidance}
\section{Laboratory 4: operational recovery narrative}
\begin{Worked}\textbf{Component}: \texttt{VC-OPS-01/1}, governed by \texttt{LG-OPS-NARR/1} from Document 03. \textbf{View}: \texttt{VW-OPS-01/1}.\par
\textbf{Trigger}: transaction remains pending beyond the operational alert threshold \texttt{[TBD]}. \textbf{First response}: verify worker queue depth and provider health; identify duplicate-safe retry rules. \textbf{Escalation owner}: \texttt{[TBD]}. \textbf{Stop criterion}: consistent status is observable and reconciled to authoritative payment records. \textbf{Known issue}: alert threshold unapproved; \texttt{INC-OPS-01}.\end{Worked}
\section{Bridging C4 descriptions to views and model kinds}
\begin{longtable}{p{0.24\linewidth}p{0.34\linewidth}p{0.28\linewidth}}\toprule\textbf{C4 resource}&\textbf{Proposed AD use}&\textbf{What must be added}\\\midrule\endhead
System context diagram&Component in a context/integration architecture view&Viewpoint, EoI, concerns, perspective mapping, version and legend/model-kind conventions.\\
Container diagram&Component in deployment or structural architecture view&Define ``container'', environment, service identity and allocations; record any cross-view relation.\\
Component diagram&Component in logical design architecture view&Specify component meaning, interface semantics, scope and governing conventions.\\
Dynamic diagram&Component in behavioural architecture view&Scenario purpose, error cases, message legend and linkage to contracts.\\
Deployment diagram&Component in allocation architecture view&Environment, runtime nodes, failure domains, deployed artifact mapping.\\\bottomrule\end{longtable}
\begin{Guidance}C4 terminology is not automatically equivalent to the 42010 vocabulary. For example, a C4 \emph{component} is an architectural building block in a software design; a 42010 \emph{architecture view component} is an identifiable portion of an AD. Record a deliberate mapping, not a one-to-one substitution.\end{Guidance}
\section{Automation and repository structure}
\begin{lstlisting}[caption={Illustrative architecture documentation layout}]
architecture/
  ad/manifest.yaml              # EoI, purpose, baselines
  registry/stakeholders.yaml
  registry/perspectives.yaml
  registry/concerns.yaml
  registry/aspects.yaml
  viewpoints/                   # VP specifications
  model-kinds/                  # MK and legend specs
  views/                        # view manifests
  components/                   # PlantUML, narratives, etc.
  correspondence/methods.yaml
  correspondence/relations.yaml
  decisions/                    # essential ADRs
  issues/                       # known inconsistencies
  evidence/                     # validation results
\end{lstlisting}
\subsection{Validation stages}
\begin{enumerate}
\item \textbf{Parse and schema-check}: verify IDs, types, dates, revisions and URIs; reject unrecognized references.
\item \textbf{Perform structural trace checks}: evaluate the concern--viewpoint--view path and mandatory metadata joins.
\item \textbf{Build components}: render diagrams and test for syntactic errors, preserving source hash and renderer version.
\item \textbf{Apply correspondence methods}: run documented predicates on pinned inventory and AD baselines.
\item \textbf{Review semantics}: a human reviewer assesses whether the view actually addresses the concern and records inconsistencies or omissions.
\item \textbf{Freeze evidence}: publish hashes, build logs, reviewer and version identifiers with the approved AD revision.
\end{enumerate}
\begin{Guidance}Use CI for repeatable checks; do not confuse CI success with a conforming AD. Even a perfectly rendered diagram might omit a fundamental stakeholder, misrepresent the EoI or fail to answer an important concern. Build instructions should explicitly require pdfLaTeX only for the present suite, but this is a repository-local compilation rule and not a requirement of the standard.\end{Guidance}
\section{Security engineering evidence as view components}
Evidence from CodeQL, GitHub code scanning, secret-scanning triage, dependency inventories or DAST can \emph{inform} a security viewpoint when the AD's concerns and scope warrant it. A SARIF artifact or alert count is neither inherently an architecture view nor architecture evidence on its own. Identify source, snapshot, applicability and interpretation, and document a method by which an architectural claim can be assessed.\par
\begin{Worked}\textbf{Concern}: ``Does every internet-facing ingress route to an identified authorization boundary?'' \textbf{Potential evidence}: interface inventory and route definitions, backed by code-level validation. \textbf{View component}: access-control flow model with interface IDs and an accompanying gap list. \textbf{Correspondence method}: match ingress interfaces with recorded enforcement points and named owners. Static-analysis results are supplemental evidence, not a substitute for the model or review.\end{Worked}
\section{Laboratory acceptance criteria}
\begin{itemize}
\item Each generated figure has stable ID, revision, governing model kind and source reference.
\item Each figure belongs to a named architecture view governed by an identified architecture viewpoint.
\item Concern, stakeholder perspective and aspect coverage are explicit rather than assumed from titles.
\item Non-model narratives have legends; model-based components use documented conventions.
\item Cross-figure inconsistencies are registered, not quietly normalized.
\item Essential decisions include rationale and links to the affected AD elements.
\item Validation artifacts distinguish automated structural checks from human semantic judgment.
\end{itemize}
\LegalFooter
\end{document}
